% Method wording prepared while the maximum-context measurement is running.
% Insert measured coverage and values only after the final run is reconciled.
We measure Qwen2.5-7B-Instruct-1M~\cite{qwen2025qwen25oneM} on two
A100 PCIe 40\,GB GPUs with tensor parallelism two, BF16 weights and KV,
and batch one. The SGLang 0.5.10rc0 worker uses native Dual Chunk Attention,
with sparse acceleration disabled. An A100-compatible FlashAttention build
and FlashInfer SiLU are fixed across all paths. The input is a fixed
synthetic text stream, truncated to each total context length $C=M+64$.
Each plotted point is one timed execution after warmup, rather than an
average over independent workloads.

HBM reuse begins with the complete historical KV resident in the recovery
GPU group. DRAM reuse restores the same history from pinned host memory,
then computes the tail; H2D time is measured directly, without assumed
bandwidth or transfer--compute overlap. Recompute begins with empty KV,
rebuilds $M$ in chunks of at most 8,192 tokens, and processes the same
separate 64-token tail. This fixed boundary preserves the execution
partition used to construct the saved state. Timing ends after synchronized
model-forward completion, before vocabulary projection. Every timed trial
checks all-layer tail KV, hidden states, and the next greedy token against
the reference execution.

These are worker-level path costs for this implementation. They exclude
fault detection, queueing, backup creation, model loading, and service
handoff. Complete-history HBM residency is a reuse reference, not a
finite-capacity Peer policy; serialized DRAM is not an optimized overlapping
recovery baseline. The measurements therefore motivate preserving reusable
state without establishing end-to-end WCT gains or model quality at 1M.
